\documentclass[10pt,a4paper]{amsart}
\usepackage[margin=19mm]{geometry}
\usepackage{amsmath,amssymb,mathtools}
\usepackage[scr=rsfs,cal=euler]{mathalfa}
\usepackage{xfrac}
\usepackage[unicode,hidelinks,bookmarks=false]{hyperref}
\newcommand{\ie}{i.e.\ }
\newcommand{\cf}{cf.\ }
\newcommand{\resp}{respectively\ }
\newcommand{\Subsection}{\subsection}
\providecommand{\nobreakdash}{\nobreak\hbox{-}}
\setlength{\emergencystretch}{2em}
\multlinegap0pt
\usepackage{mathtools}
\usepackage{cancel}
\usepackage{amssymb}
\usepackage[shortlabels]{enumitem}
\usepackage{csquotes}
\usepackage{xcolor}
\usepackage{tikz}
\newtheorem{lemma}{Lemma}
\newcommand{\DIV}{\operatorname{\mathrm{div}}}
\newcommand{\Id}{\mathrm{I}_3}
\def\O{\Omega}
\def\be{\boldsymbol{\epsilon}}
\def\bsig{\boldsymbol{\sigma}}
\def\ddgs{\boldsymbol{X}_{{\bf grad},h,0}}
\newcommand{\dgr}[1]{\mathbb G^1_{#1}}
\def\dgs{\boldsymbol{X}_{{\bf grad},h}}
\newcommand{\norm}[2]{\|#2\|_{#1}}
\def\tr{{\rm {tr}}}
\newcommand{\uvec}[1]{\underline{\mathbf{#1}}}
\title{A higher order polyhedral method for contact mechanics with Tresca friction}
\author{Jerome Droniou, Raman Kumar, Roland Masson,, Ritesh Singla}
\date{Source: arXiv:2601.07586v2 (CC BY 4.0)}
\begin{document}
\maketitle

\noindent\emph{Source and licence.} A higher order polyhedral method for contact mechanics with Tresca friction. Version arXiv:2601.07586v2 (CC BY 4.0), distributed by arXiv under the Creative Commons Attribution 4.0 International licence, http://creativecommons.org/licenses/by/4.0/, as recorded in the arXiv licence metadata for this exact version and shown on the abstract page https://arxiv.org/abs/2601.07586v2. Full licence notice: \texttt{licenses/arxiv-2601.07586-CC-BY-4.0.txt}.\par\smallskip

\noindent\textit{Editorial scope.} Counted unit: the article's lemma eq:bound.sigma.h (source0.tex lines 850--857). The article's complete original proof of it is delivered with it (source0.tex lines 859--884, one \texttt{proof} environment of 26 lines). No mathematics is written, repaired or re-derived: the statement and the proof are the article's own lines, in the article's own notation, and no retained line is substituted or re-flowed.  Retained uncounted as the source-local context the proof needs: lines 423--428; lines 468--470.  Nothing was omitted from any retained span, so the package carries no omission record.  No retained line contains a citation, so the wrapper carries no bibliography and no bibliography entry.  No retained line contains a figure environment, an image-inclusion command or drawing code, so no image is delivered and the package declares no asset dependency.  Editorial formatting only, in addition: an AMS \texttt{amsart} preamble that restates the article's own declarations and macros used by the retained lines, namely the environments \texttt{lemma} on separate counters, together with the standard packages those lines need.  The retained environments print in this document as Lemma 1 in order of appearance, whereas the article prints the same items with the numbering of its own full document.  The complete native source of the article as distributed in the arXiv e-print for this exact version is bundled unchanged as \texttt{sources/192515/source0.tex}.
Define the discrete symmetric gradient $\epsilon_h$, divergence $\DIV_h$, and stress tensor $\sigma_h$ as: for all $\uvec{v}_h\in\dgs$,
\begin{alignat}{2}
\be_h(\uvec{v}_h)&\coloneqq\frac{1}{2}(\dgr{h}\uvec{v}_h+(\dgr{h}\uvec{v}_h)^T),\quad\DIV_h(\uvec{v}_h)\coloneqq \tr(\be_h(\uvec{v}_h)),\nonumber\\
 \bsig_h(\uvec{v}_h)&\coloneqq \frac{E}{1+\nu}\left(\be_h(\uvec{v}_h)+\frac{\nu}{1-2\nu} \DIV_h(\uvec{v}_h)\Id\right).
 \label{eq:def.sigma.h}
\end{alignat}

\begin{align}\label{eq:def.inner.Eh}
\langle\uvec{v}_h,\uvec{w}_h\rangle_{\mathrm{en},h}\coloneqq\int_{\Omega}\bsig_h(\underline{\bf v}_h)\colon\be_h(\underline{\bf w}_h)+\mu_{1}\mathbb S_h(\underline{\bf v}_h,\underline{\bf w}_h),
\end{align}

\begin{lemma}[Bound on the discrete stress]%\label{lem:stress_bdness}
For all $\uvec{v}_{h}\in\ddgs$, it holds: 
\begin{equation}\label{eq:bound.sigma.h}
\norm{L^2(\O)}{\bsig_h(\uvec{v}_h)}\lesssim \sqrt{\mu_{2}} \left(1+\left(\frac{\nu}{1-2\nu}\right)^{\frac{1}{2}} \right) \norm{\mathrm{en},h}{\uvec{v}_h},
%\label{lem:stress_bdness_1}
\end{equation}
where the hidden constant in $\lesssim$ is independent of $E$ and $\nu.$
\end{lemma}

\begin{proof}
First, recall the definition of $\bsig_h(\cdot)$ from \eqref{eq:def.sigma.h} to get
\begin{equation*}
    \bsig_h(\uvec{v}_h)= \frac{E}{1+\nu}\left(\be_h(\uvec{v}_h)+\frac{\nu}{1-2\nu} \DIV_h(\uvec{v}_h)\Id\right).
\end{equation*}
An application of the triangle inequality in the above equation leads to
\begin{equation}
\norm{L^2(\O)}{\bsig_h(\uvec{v}_h)} \lesssim \frac{E}{1+\nu}\left(\norm{L^2(\O)}{\be_h(\uvec{v}_h)}+\frac{\nu}{1-2\nu}\norm{L^2(\O)}{ \DIV_h(\uvec{v}_h)}\right). \label{l2_bd1}
\end{equation}
The definition of $\norm{\mathrm{en},h}{{\cdot}}$ from \eqref{eq:def.inner.Eh} and using $(a+b)^{1/2}\ge \frac12 (a^{1/2}+b^{1/2})$, we have
$$
\norm{\mathrm{en},h}{\uvec{v}_h}\ge \frac12 \left(\frac{E}{1+\nu}\right)^{\frac{1}{2}}
\left(\norm{L^2(\O)}{\be_h(\uvec{v}_h)}+\left(\frac{\nu}{1-2\nu}\right)^{\frac{1}{2}}\norm{L^2(\O)}{\DIV_h(\uvec{v}_h)}\right),
$$
which gives
$$
\norm{L^2(\O)}{\be_h(\uvec{v}_h)}+\left(\frac{\nu}{1-2\nu}\right)^{\frac{1}{2}}\norm{L^2(\O)}{\DIV_h(\uvec{v}_h)}\lesssim
\left(\frac{E}{1+\nu}\right)^{-\frac{1}{2}}\norm{\mathrm{en},h}{\uvec{v}_h}.
$$
Plugging this relation into \eqref{l2_bd1} leads to
$$
\norm{L^2(\O)}{\bsig_h(\uvec{v}_h)}
\lesssim \left(\frac{E}{1+\nu}\right)^{\frac{1}{2}}\left(1+\left(\frac{\nu}{1-2\nu}\right)^{\frac{1}{2}}\right)\norm{\mathrm{en},h}{\uvec{v}_h},
$$
The above inequality with $\frac{E}{1+\nu} \leq \mu_{2}$ implies \eqref{eq:bound.sigma.h}. 
\end{proof}

\end{document}
